\clearpage
\section{Прогностическая полезность характеристик команды}
\label{sec:predictors}

В этом разделе представлены результаты первого прогнозного эксперимента: сравнение сведений о наборе и возрасте компании с добавлением размера команды и числа основателей. Исследуются 3\,046 компаний YC, активных 13.07.2023, с известным статусом на 01.08.2026; положительный класс содержит 272 поглощённые или публичные компании. Все оценки получены по внешним проверочным прогнозам в соответствии с процедурой раздела~\ref{sec:modeling-methods}. Они характеризуют полезность выбранных признаков для данной задачи, а не причины успеха предпринимательских инициатив.

\subsection{Качество прогноза в основной схеме}

В табл.~\ref{tab:models-main} приведены метрики, усреднённые по двум стратифицированным повторам. Постоянный прогноз даёт логарифмическую потерю 0,3009. Использование только набора в логистической регрессии почти не меняет её: 0,3007. Добавление возраста снижает потерю до 0,2996, однако последующее включение двух характеристик команды увеличивает её до 0,3001. Следовательно, более подробное описание компании не приводит автоматически к лучшему вероятностному прогнозу.

\begin{table}[!htbp]
    \centering
    \small
    \caption{Качество внешних проверочных прогнозов для 3\,046 компаний YC с известным исходом. Признаки относятся к 13.07.2023, статусы --- к 01.08.2026. Представлены средние метрики двух стратифицированных повторов вложенной проверки. Для LogLoss и ошибки Брайера меньшее значение лучше; для AP и ROC-AUC --- большее. ЛР обозначает логистическую регрессию.}
    \label{tab:models-main}
    \begin{tabular}{@{}>{\raggedright\arraybackslash}p{7.1cm}rrrr@{}}
\toprule
Модель & LogLoss & AP & ROC-AUC & Брайер \\
\midrule
Постоянный прогноз & 0,3009 & 0,0893 & 0,5000 & 0,0813 \\
ЛР: набор и возраст & 0,3002 & 0,1028 & 0,5476 & 0,0813 \\
ЛР: набор, возраст и команда & 0,3005 & 0,1065 & 0,5388 & 0,0813 \\
ЛР: команда и география & 0,2997 & 0,1125 & 0,5659 & 0,0812 \\
ЛР: команда и теги & 0,2953 & 0,1398 & 0,6122 & 0,0804 \\
ЛР: команда и краткое описание & 0,2996 & 0,1102 & 0,5504 & 0,0812 \\
ЛР: все блоки признаков & 0,2944 & 0,1436 & 0,6193 & 0,0803 \\
CatBoost: набор, возраст и команда & 0,2997 & 0,1163 & 0,5500 & 0,0812 \\
\bottomrule
\end{tabular}

\end{table}

У CatBoost модель по набору и возрасту имеет потерю 0,2994, а модель с двумя характеристиками команды --- 0,2998. Наименьшая средняя потеря среди рассмотренных вариантов равна 0,2992 и наблюдается у CatBoost без числа основателей, то есть с набором, возрастом и размером команды. Её отличие от постоянного прогноза составляет около 0,0017. Это лучшая точечная оценка в проведённом сравнении, но не доказательство превосходства такой конфигурации на новых данных.

Способность различать положительные и отрицательные исходы остаётся ограниченной. Наибольшая средняя ROC-AUC равна 0,5622, AP --- 0,1095; обе оценки также относятся к CatBoost без числа основателей. Для модели с обоими признаками команды ROC-AUC составляет 0,5472, AP --- 0,1040. Эти результаты не дают основания представлять модель как инструмент надёжного отбора будущих поглощённых или публичных компаний. В частности, AP не означает долю правильных ответов при определённом пороге вероятности.

Постоянный прогноз в таблице не имеет ROC-AUC строго 0,5, поскольку его вероятность одинакова внутри одной проверочной части, но определяется долей исходов в соответствующей обучающей части. При объединении внешних прогнозов эти вероятности немного различаются. Это техническое следствие процедуры оценки, а не использование моделью индивидуальных характеристик компании. Значения AP такого прогноза по той же причине не обязаны точно совпадать с общей долей положительного класса.

\subsection{Дополнительная информация о команде}

Парные сравнения в табл.~\ref{tab:model-gains} и на рис.~\ref{fig:model-gains} показывают изменение логарифмической потери на одинаковых проверочных компаниях. Знак определён так, что положительная разность означает улучшение новой модели. Добавление возраста к набору в логистической регрессии даёт среднее улучшение 0,00108. Добавление команды к набору и возрасту, напротив, ухудшает среднюю оценку на 0,00045 для логистической регрессии и на 0,00046 для CatBoost.

\begin{table}[!htbp]
    \centering
    \small
    \caption{Парные изменения логарифмической потери в основной проверке по 3\,046 компаниям YC. Положительное значение означает уменьшение ошибки. Границы получены повторным выбором компаний при фиксированных прогнозах и не являются полной оценкой неопределённости процедуры обучения.}
    \label{tab:model-gains}
    \begin{tabular}{@{}>{\raggedright\arraybackslash}p{6.0cm}rrr@{}}
\toprule
Сравнение & $\Delta L$ & Компании & Наборы \\
\midrule
Характеристики команды & -0,00035 & [-0,00099; 0,00029] & [-0,00116; 0,00057] \\
География & 0,00060 & [-0,00099; 0,00214] & [-0,00142; 0,00232] \\
Тематические теги & 0,00510 & [0,00179; 0,00851] & [0,00163; 0,00808] \\
Краткое описание & 0,00051 & [-0,00035; 0,00144] & [-0,00039; 0,00110] \\
География, теги и описание совместно & 0,00609 & [0,00251; 0,00984] & [0,00204; 0,00964] \\
CatBoost против ЛР на одинаковых признаках & 0,00146 & [-0,00054; 0,00331] & [-0,00025; 0,00331] \\
Полная модель против постоянного прогноза & 0,00525 & [0,00104; 0,00955] & [0,00130; 0,00954] \\
Полная модель против модели с тегами & 0,00099 & [-0,00024; 0,00223] & [-0,00008; 0,00197] \\
Исключение тегов из полной модели & 0,00450 & [0,00147; 0,00763] & [0,00135; 0,00723] \\
Исключение описания из полной модели & 0,00051 & [-0,00010; 0,00115] & [0,00008; 0,00086] \\
Исключение страны из полной модели & 0,00040 & [-0,00063; 0,00138] & [-0,00082; 0,00148] \\
Исключение команды из полной модели & -0,00023 & [-0,00115; 0,00066] & [-0,00123; 0,00082] \\
\bottomrule
\end{tabular}

\end{table}

\begin{figure}[!htbp]
    \centering
    \includegraphics[width=\textwidth]{modeling/02_incremental_value.pdf}
    \caption{Изменение проверочной логарифмической потери при добавлении признаков и замене алгоритма для компаний YC, активных 13.07.2023, с известным статусом на 01.08.2026 ($n=3\,046$). Точка --- средняя парная разность двух стратифицированных повторов; значения справа от нуля означают улучшение. Отрезки охватывают центральные 95\% из 2\,000 повторных выборок компаний при фиксированных прогнозах, без переобучения моделей и учёта зависимости внутри наборов.}
    \label{fig:model-gains}
\end{figure}

Исключение признаков по отдельности уточняет результат для CatBoost. Добавление размера команды при уже известном числе основателей даёт практически нулевую среднюю разность: 0,00004. Добавление числа основателей при известном размере команды ухудшает среднюю потерю на 0,00064. Следовательно, в рассмотренных конфигурациях число основателей не даёт убедительного дополнительного улучшения; однако это не доказывает его бесполезности при других представлениях, алгоритмах или исходах.

Во всех шести сравнениях основной схемы условные границы охватывают ноль. С учётом ограничений их расчёта это не заменяет формальной проверки гипотез, но показывает, что небольшие средние различия не следует интерпретировать как точный рейтинг признаков. Неустойчивость заметна и между двумя повторами: добавление команды в CatBoost увеличивает потерю примерно на 0,00129 в первом повторе и уменьшает на 0,00037 во втором. Среднее скрывает смену направления различия.

Таким образом, наблюдавшиеся в EDA различия размера команды между группами позднего статуса не превратились в устойчивое улучшение прогноза после учёта возраста и набора. Групповое различие и дополнительная прогностическая информация являются разными результатами исследования.

\subsection{Проверка с разделением наборов YC}

В дополнительной схеме компании одного набора не встречаются одновременно в обучении и проверке. Результаты приведены в табл.~\ref{tab:models-grouped}; рис.~\ref{fig:model-quality} сопоставляет логарифмические потери двух схем. Внешние проверочные части содержат от 584 до 655 компаний и от 51 до 57 положительных исходов, то есть различия не сводятся к отсутствию положительного класса в какой-либо части.

\begin{table}[!htbp]
    \centering
    \small
    \caption{Качество прогнозов с разделением по полному коду набора YC: один повтор вложенной проверки, 3\,046 компаний с известным статусом на 01.08.2026 и исходными признаками на 13.07.2023. Каждый набор целиком отнесён к одной внешней части; внутри обучения разделение по наборам сохранено. Направления улучшения метрик совпадают с табл.~\ref{tab:models-main}.}
    \label{tab:models-grouped}
    \begin{tabular}{@{}>{\raggedright\arraybackslash}p{7.1cm}rrrr@{}}
\toprule
Модель & LogLoss & AP & ROC-AUC & Брайер \\
\midrule
Постоянный прогноз & 0,3011 & 0,0893 & 0,5000 & 0,0813 \\
ЛР: набор и возраст & 0,3016 & 0,1018 & 0,5353 & 0,0816 \\
ЛР: набор, возраст и команда & 0,3019 & 0,1017 & 0,5220 & 0,0816 \\
ЛР: команда и география & 0,3013 & 0,1087 & 0,5549 & 0,0815 \\
ЛР: команда и теги & 0,2968 & 0,1353 & 0,6052 & 0,0808 \\
ЛР: команда и краткое описание & 0,3014 & 0,1037 & 0,5396 & 0,0816 \\
ЛР: все блоки признаков & 0,2958 & 0,1402 & 0,6154 & 0,0806 \\
CatBoost: набор, возраст и команда & 0,3004 & 0,1109 & 0,5286 & 0,0813 \\
\bottomrule
\end{tabular}

\end{table}

\begin{figure}[!htbp]
    \centering
    \includegraphics[width=\textwidth]{modeling/01_model_quality.pdf}
    \caption{Логарифмическая потеря восьми вариантов прогноза статуса на 01.08.2026 по характеристикам компаний YC на 13.07.2023 ($n=3\,046$ с известным исходом). Круги --- среднее двух стратифицированных повторов, квадраты --- проверка с разделением наборов YC. Меньшее значение означает лучший вероятностный прогноз. Все точки получены по внешним проверочным наблюдениям; разделение наборов не является проверкой нового календарного периода.}
    \label{fig:model-quality}
\end{figure}

При разделении наборов постоянный прогноз имеет потерю 0,3011. Логистическая регрессия с возрастом и командой даёт 0,3021, то есть оказывается хуже этого ориентира. CatBoost с теми же признаками имеет потерю 0,2998 и превосходит соответствующую логистическую регрессию примерно на 0,0023. Однако относительно CatBoost только по набору и возрасту выигрыш от добавления команды составляет лишь 0,00021; в основной схеме знак этого сравнения был противоположным.

Ранжирование также ухудшается относительно основной схемы: у CatBoost с двумя характеристиками команды ROC-AUC составляет 0,5294, AP --- 0,0966. У варианта без числа основателей эти показатели равны 0,5320 и 0,0993. Это указывает на ограниченную переносимость обнаруженных зависимостей между наборами в рамках проведённой проверки. Одного группового разбиения недостаточно для общей оценки всех возможных различий между наборами, а тем более для вывода о качестве прогнозирования после августа 2026 года.

\subsection{Вероятности и наблюдаемые частоты}

Рисунок~\ref{fig:model-calibration} сопоставляет средние проверочные вероятности с фактическими долями положительного класса в группах примерно одинакового размера. Для всех трёх моделей использован первый стратифицированный повтор, выбранный до сравнения результатов. Границы групп определяются отдельно по прогнозам каждой модели; принадлежность компании к группам разных моделей поэтому может различаться.

\begin{figure}[!htbp]
    \centering
    \includegraphics[width=\textwidth]{modeling/03_calibration.pdf}
    \caption{Средние вероятности и наблюдаемые доли поглощённых или публичных компаний на 01.08.2026 среди исходно активных компаний YC ($n=3\,046$). Показаны внешние проверочные прогнозы первого стратифицированного повтора по признакам на 13.07.2023. Каждая точка объединяет примерно пятую часть компаний по предсказанной вероятности конкретной модели. Пунктир соответствует совпадению средней вероятности и наблюдаемой частоты; дополнительная калибровка не выполнялась.}
    \label{fig:model-calibration}
\end{figure}

Для CatBoost с возрастом и командой средняя предсказанная вероятность в нижней группе составляет 6,00\%, тогда как наблюдаемая доля положительного исхода --- 7,52\%. В верхней группе эти значения равны 11,31\% и 9,69\%. В данном повторе модель раздвигает средние вероятности крайних групп сильнее, чем различаются наблюдаемые частоты. При этом наибольшая фактическая доля наблюдается не в верхней группе вероятности, а в предпоследней: 10,20\%.

У логистической регрессии с теми же признаками крайние группы имеют средние вероятности 6,48\% и 11,35\% при наблюдаемых долях 6,39\% и 10,84\%. Такое совпадение в отдельных группах не означает достаточного разделения исходов: большинство предсказаний остаётся близко к общей частоте положительного класса. Кроме того, небольшие различия групповых частот подвержены случайным колебаниям. По пяти точкам одного повтора нельзя установить систематическую ошибку калибровки либо выбрать способ её исправления без дополнительной проверки.

\subsection{Вклад признаков в прогнозы CatBoost}

На рис.~\ref{fig:model-shap} представлены средние модули SHAP для CatBoost с пятью исходными признаками. Каждая компания объясняется моделью, для которой она была внешним проверочным наблюдением. Поэтому рисунок не основан на объяснении обучающих прогнозов, но остаётся описанием поведения уже выбранного семейства моделей.

\begin{figure}[!htbp]
    \centering
    \includegraphics[width=\textwidth]{modeling/04_shap_importance.pdf}
    \caption{Средние абсолютные вклады SHAP в проверочные прогнозы CatBoost для 3\,046 компаний YC с известным статусом на 01.08.2026. Используются признаки на 13.07.2023 и пять внешних моделей первого стратифицированного повтора. Единица вклада --- логарифм шансов положительного статуса, а не процент вероятности. Длина столбца не показывает направление связи, причинный эффект или улучшение качества после добавления признака.}
    \label{fig:model-shap}
\end{figure}

Наибольший средний модуль вклада имеет логарифм размера команды --- 0,0831. Далее следуют год набора (0,0719), возраст (0,0656), число основателей (0,0521) и сезон набора (0,0315). Эти числа характеризуют изменение исходного выхода модели относительно её базового значения при разложении SHAP. Они не являются коэффициентами логистической регрессии и не означают, например, увеличение вероятности на 8,31 процентного пункта из-за размера команды.

Результат особенно важен в сопоставлении с исключением признаков: размер команды имеет наибольший средний вклад в прогнозы, однако его добавление при известном числе основателей почти не меняет среднюю проверочную потерю. Признак может заметно влиять на ответы модели, не обеспечивая устойчивого улучшения этих ответов относительно модели без него. Возможны дублирование информации другими признаками и использование зависимостей, недостаточно устойчивых вне обучения; проведённое сравнение не позволяет однозначно выбрать одно объяснение.

Следовательно, рейтинг SHAP не заменяет проверку дополнительной полезности. На этом этапе нельзя называть размер команды главным доказанным предиктором успеха только потому, что его столбец самый длинный. Корректный вывод уже: в рассмотренном CatBoost этот признак сильнее остальных в среднем участвует в формировании проверочных прогнозов.

\subsection{Выводы первого эксперимента}

Проверка не выявила устойчивого улучшения вероятностного прогноза от совместного добавления размера команды и числа основателей к набору и возрасту. Между алгоритмами и схемами оценки есть различия, но абсолютное качество ранжирования остаётся невысоким. Лучшие точечные оценки принадлежат CatBoost без числа основателей; считать этот вариант окончательно выбранной моделью на основании текущих результатов преждевременно.

Результат уточняет смысл описательных зависимостей: различия между группами последующего статуса не гарантируют полезности для прогноза отдельной компании. Он также показывает, почему объяснение модели через SHAP должно сопровождаться сравнением проверочного качества.

Вывод ограничен выбранным показателем, когортой с известным исходом, пятью характеристиками и заданными конфигурациями моделей. Пока не проверены география, тематические теги, текстовое представление деятельности, альтернативная обработка сомнительных нулевых значений и влияние семи публичных компаний на объединённый положительный класс. Их исследование должно продолжать зафиксированную процедуру сравнения, а не служить поиском отдельного удачного результата. Новое изменение модели после просмотра этих оценок потребует явно учитывать исследовательский характер такого выбора.
\clearpage
